\documentclass[11pt,a4paper]{article}

\usepackage{main-sss}
\usepackage{corrige-sss}

\renewcommand{\matiere}{Physique-Chimie}
\renewcommand{\examen}{DNB}
\renewcommand{\annee}{2026}
\renewcommand{\localisation}{Pays du Groupe 1}
\renewcommand{\titresujet}{La grêle}
\renewcommand{\niveau}{Brevet}
\renewcommand{\typedocument}{Corrigé}
\renewcommand{\identifiantsujet}{26GENSCG11}

\begin{document}

\pagination
\titre

\q[1 point]
Les orages de grêles peuvent avoir des conséquences matérielles (sur les cultures, les véhicules, les bâtiments) et sur les êtres vivants (personnes, animaux).

\q[0,5 point]
Le grêlon est à l'état solide. La modélisation correspondante est la B.

\q[1 point]
La molécule d'eau contient 2 atomes d'hydrogène (H) et 1 atome d'oxygène (O).

\q[1 point]
La formation de grêlons correspond à un changement d'état de l'eau liquide vers l'eau solide. Il s'agit donc d'une transformation physique (les espèces chimiques restent les mêmes).

\q[2 point]
Durant les 5 premières secondes de la chute, les grêlons passent d'une vitesse nulle à une vitesse de \SI{25}{m/s}. Ils ont donc un mouvement accéléré.

De plus, comme précisé dans le texte, ils chutent verticalement, leur trajectoire est donc une droite et leur mouvement est rectiligne.

\q[1 point]
Au cours de la chute du grêlon, son énergie potentielle de pesanteur est convertie en énergie cinétique.

\q[1 point]
La formule de l'énergie cinétique est donnée par la relation B : $E_c = \frac{1}{2} \times m \times v^2$.

\q[2,5 points]
La vitesse maximale atteinte par le grêlon est de \SI{25}{m/s}. Son énergie cinétique vaut alors :

$E_c = \frac{1}{2} \times m \times v^2 = \frac{1}{2} \times 0,013 \times 25^2 = 4,1 J $

D'après le document 3, une énergie cinétique de 4,1 J correspond à des grêlons d'un diamètre de 25 à 40 mm. Ils peuvent causer de graves dommages aux cultures et endommager les voitures.


\end{document}